\documentclass{article}
\usepackage[utf8]{inputenc}
\usepackage[T1]{fontenc}
\usepackage{amsmath}
\usepackage{amsfonts}
\usepackage{amssymb}
\usepackage{enumitem}

\usepackage[margin=0.7in]{geometry}

\usepackage{tikz}
\usetikzlibrary{matrix,arrows,decorations.pathmorphing,shapes.geometric,calc,snakes,backgrounds,arrows.meta}
\usepackage{xcolor}

\begin{document}
\pagestyle{plain}

\section*{Solutions: Discrete Mathematics, IMC 2001--2006}

\begin{center}
April 2026
\end{center}

\subsection*{Combinatorics}

\textbf{Problem 1.} (IMC 2001, Day 1, Problem 1)
Let \( n \) be a positive integer. Consider an \( n \times n \) matrix with entries \( 1, 2, \ldots, n^2 \) written in order starting from the top-left corner and moving along each row in turn left-to-right. We choose \( n \) entries of the matrix such that exactly one entry is chosen in each row and each column. What are the possible values of the sum of the selected entries?

\textbf{Solution.}
The entry in row \( j \), column \( k \) equals \( n(j-1) + k \). The choice corresponds to a permutation \( \sigma \in S_n \), selecting entries \( \{(j, \sigma(j)) : j = 1, \ldots, n\} \). The sum is
\[
\sum_{j=1}^{n} \bigl(n(j-1) + \sigma(j)\bigr) = n \cdot \frac{n(n-1)}{2} + \sum_{j=1}^{n} \sigma(j) = \frac{n(n-1) \cdot n}{2} + \frac{n(n+1)}{2} = \frac{n(n^2+1)}{2},
\]
which is independent of \( \sigma \). The only possible value is \( \dfrac{n(n^2+1)}{2} \). \hfill\(\square\)
\\

\textbf{Problem 2.} (IMC 2002, Day 2, Problem 2)
Two hundred students participated in a mathematical contest. They had 6 problems to solve. It is known that each problem was correctly solved by at least 120 participants. Prove that there must be two participants such that every problem was solved by at least one of these two students.

\textbf{Solution.}
For each pair of students, consider the set of problems not solved by either of them. There are \( \binom{200}{2} = 19900 \) pairs; we must prove that at least one such set is empty.

For each problem, at most 80 students did not solve it, contributing at most \( \binom{80}{2} = 3160 \) pairs. Over all 6 problems, at most \( 6 \cdot 3160 = 18960 \) pairs have a non-empty set.

Since \( 19900 - 18960 = 940 > 0 \), at least 940 pairs have an empty set. For any such pair, every problem was solved by at least one of the two students. \hfill\(\square\)
\\

\textbf{Problem 3.} (IMC 2003, Day 2, Problem 4)
Find all positive integers \( n \) for which there exists a family \( \mathcal{F} \) of three-element subsets of \( S = \{1, 2, \ldots, n\} \) satisfying:
\begin{enumerate}
    \item for any two different elements \( a, b \in S \), there exists exactly one \( A \in \mathcal{F} \) containing both \( a \) and \( b \);
    \item if \( a, b, c, x, y, z \) are elements of \( S \) such that \( \{a, b, x\}, \{a, c, y\}, \{b, c, z\} \in \mathcal{F} \), then \( \{x, y, z\} \in \mathcal{F} \).
\end{enumerate}

\textbf{Solution.}
\textit{Answer:} \( n = 2^m - 1 \) for some positive integer \( m \).

Condition (i) defines an operation \( * \) on \( S \) by \( a * b = c \) iff \( \{a,b,c\} \in \mathcal{F} \) (for \( a \neq b \)). This operation satisfies: (a) \( a \neq a*b \neq b \); (b) \( a*b = b*a \); (c) \( a*(a*b) = b \).

Condition (ii) gives a form of associativity: \( x*(a*c) = (x*a)*c \) when the arguments are distinct.

Add an element 0 to \( S \) and define \( a*a = 0 \) and \( a*0 = a \). Then \( (S \cup \{0\}, *) \) becomes a finite abelian group in which every non-identity element has order 2. Such a group must be isomorphic to \( (\mathbb{Z}/2\mathbb{Z})^m \), hence has order \( 2^m \), giving \( n + 1 = 2^m \).

Conversely, for \( n = 2^m - 1 \), identify \( S \) with \( (\mathbb{Z}/2\mathbb{Z})^m \setminus \{0\} \) and define \( a * b \) as the componentwise XOR (addition mod 2). This satisfies both conditions. \hfill\(\square\)
\\

\textbf{Problem 4.} (IMC 2004, Day 1, Problem 1)
Let \( S \) be an infinite set of real numbers such that \( |s_1 + s_2 + \cdots + s_k| < 1 \) for every finite subset \( \{s_1, s_2, \ldots, s_k\} \subset S \). Show that \( S \) is countable.

\textbf{Solution.}
For each positive integer \( n \), let \( S_n^+ = S \cap (1/n, \infty) \) and \( S_n^- = S \cap (-\infty, -1/n) \).

If \( |S_n^+| \geq n \), then choosing \( n \) elements from \( S_n^+ \) gives a sum exceeding \( n \cdot (1/n) = 1 \), contradicting the hypothesis. So \( |S_n^+| < n \), and similarly \( |S_n^-| < n \).

Every nonzero element of \( S \) belongs to some \( S_n^+ \) or \( S_n^- \), so
\[
S \subset \{0\} \cup \bigcup_{n=1}^{\infty} (S_n^+ \cup S_n^-),
\]
a countable union of finite sets, hence countable. \hfill\(\square\)
\\

\textbf{Problem 5.} (IMC 2005, Day 1, Problem 2)
For an integer \( n \geq 3 \) consider the sets
\begin{align*}
S_n &= \{(x_1, \ldots, x_n) : \forall\, i,\; x_i \in \{0,1,2\}\}, \\
A_n &= \{(x_1, \ldots, x_n) \in S_n : \forall\, i \leq n-2,\; |\{x_i, x_{i+1}, x_{i+2}\}| \neq 1\}, \\
B_n &= \{(x_1, \ldots, x_n) \in S_n : \forall\, i \leq n-1,\; (x_i = x_{i+1} \Rightarrow x_i \neq 0)\}.
\end{align*}
Prove that \( |A_{n+1}| = 3 \cdot |B_n| \).

\textbf{Solution.}
Let \( a_n = |A_n| \) and \( b_n = |B_n| \). We show both satisfy the same recurrence.

Define \( A'_n = \{x \in A_n : x_{n-1} = x_n\} \) and \( A''_n = A_n \setminus A'_n \). Similarly, \( B'_n = \{x \in B_n : x_n = 0\} \) and \( B''_n = B_n \setminus B'_n \).

For the \( A \)-sequences: an element of \( A'_{n+1} \) ends with two equal entries, so the preceding entry must differ (by the ``no three equal'' condition), hence \( |A'_{n+1}| = |A''_n| \). For \( A''_{n+1} \): the last entry differs from the second-to-last, and there are 2 choices for it regardless of whether the penultimate pair was equal or not, giving \( |A''_{n+1}| = 2|A'_n| + 2|A''_n| \). Thus
\[
a_{n+1} = |A'_{n+1}| + |A''_{n+1}| = |A''_n| + 2|A'_n| + 2|A''_n| = 2a_n + |A''_n| - |A'_n| + 2|A'_n|.
\]
More directly: \( a_{n+1} = 2a_n + a'_n \) with \( a'_{n+1} = a''_n = a_n - a'_n \), which yields \( a_{n+2} = 2a_{n+1} + 2a_n \).

For the \( B \)-sequences: similarly \( |B'_{n+1}| = |B''_n| \) and \( |B''_{n+1}| = 2|B'_n| + 2|B''_n| \), giving \( b_{n+2} = 2b_{n+1} + 2b_n \).

Initial values: \( a_3 = 24 \), \( a_4 = 72 \) and \( b_2 = 8 \), \( b_3 = 24 \). We verify \( a_4 = 3 \cdot b_3 = 72 \) and \( a_3 = 3 \cdot b_2 = 24 \). Since both sequences satisfy the same recurrence with \( a_{n+1} = 3 b_n \) for two consecutive initial values, by induction \( |A_{n+1}| = 3 \cdot |B_n| \) for all \( n \geq 2 \). \hfill\(\square\)

\newpage

\subsection*{Number Theory and Algebra}

\textbf{Problem 6.} (IMC 2001, Day 1, Problem 2)
Let \( r, s, t \) be positive integers which are pairwise relatively prime. If \( a \) and \( b \) are elements of a commutative multiplicative group with unity element \( e \), and \( a^r = b^s = (ab)^t = e \), prove that \( a = b = e \). Does the same conclusion hold if \( a \) and \( b \) are elements of an arbitrary non-commutative group?

\textbf{Solution.}
Since \( \gcd(s, t) = 1 \), there exist integers \( u, v \) with \( us + vt = 1 \). Using commutativity:
\[
ab = (ab)^{us+vt} = (ab)^{us} \cdot ((ab)^t)^v = a^{us} b^{us} = a^{us} (b^s)^u = a^{us}.
\]
Therefore \( b = a^{us - 1} \), so \( b^r = a^{r(us-1)} = (a^r)^{us-1} = e \). Now \( b^r = b^s = e \) and \( \gcd(r, s) = 1 \), so \( b = b^{xr + ys} = e \) for suitable integers \( x, y \). Similarly \( a = e \).

The conclusion fails for non-commutative groups. In \( S_7 \), take \( a = (1\;2\;3) \) and \( b = (3\;4\;5\;6\;7) \). Then \( a^3 = b^5 = e \) and \( ab = (1\;2\;3\;4\;5\;6\;7) \), so \( (ab)^7 = e \). Here \( 3, 5, 7 \) are pairwise coprime, yet \( a \neq e \neq b \). \hfill\(\square\)
\\

\textbf{Problem 7.} (IMC 2002, Day 2, Problem 3)
For each \( n \geq 1 \) let
\[
a_n = \sum_{k=0}^{\infty} \frac{k^n}{k!}, \qquad b_n = \sum_{k=0}^{\infty} (-1)^k \frac{k^n}{k!}.
\]
Show that \( a_n \cdot b_n \) is an integer.

\textbf{Solution.}
We prove that \( a_n / e \) and \( b_n \cdot e \) are integers for all \( n \geq 0 \) (with \( 0^0 = 1 \)).

Base case: \( a_0 = \sum_{k=0}^{\infty} \frac{1}{k!} = e \) and \( b_0 = \sum_{k=0}^{\infty} \frac{(-1)^k}{k!} = e^{-1} \). So \( a_0/e = 1 \) and \( b_0 \cdot e = 1 \).

Inductive step: using \( (k+1)^n = \sum_{m=0}^{n} \binom{n}{m} k^m \) and reindexing:
\[
a_{n+1} = \sum_{k=0}^{\infty} \frac{(k+1)^n}{k!} = \sum_{m=0}^{n} \binom{n}{m} a_m,
\]
and similarly \( b_{n+1} = -\sum_{m=0}^{n} \binom{n}{m} b_m \). (The sign change comes from the \( (-1)^k \) factor.)

If \( a_0, \ldots, a_n \) are integer multiples of \( e \), then so is \( a_{n+1} \). Similarly for \( b \). Hence \( a_n \cdot b_n = \frac{a_n}{e} \cdot (b_n \cdot e) \) is a product of two integers. \hfill\(\square\)
\\

\textbf{Problem 8.} (IMC 2003, Day 1, Problem 2)
Let \( a_1, a_2, \ldots, a_{51} \) be non-zero elements of a field. We simultaneously replace each element with the sum of the 50 remaining ones, obtaining \( b_1, \ldots, b_{51} \). If this new sequence is a permutation of the original one, what can be the characteristic of the field?

\textbf{Solution.}
\textit{Answer:} The characteristic is 2 or 7.

Let \( S = a_1 + a_2 + \cdots + a_{51} \). Then \( b_i = S - a_i \). Since \( (b_i) \) is a permutation of \( (a_i) \):
\[
\sum b_i = \sum a_i \implies 51S - S = S \implies 49S = 0.
\]

\textit{Case 1:} Characteristic does not divide 49, i.e., \( \neq 7 \). Then \( S = 0 \), so \( b_i = -a_i \). The map \( a_i \mapsto -a_i \) must be a permutation. If the characteristic is also \( \neq 2 \), then \( -a_i \neq a_i \) for all \( i \) (since \( a_i \neq 0 \)), so the permutation decomposes into pairs \( \{a_i, -a_i\} \). But 51 is odd, contradiction. So the characteristic must be 2.

\textit{Case 2:} Characteristic is 7. Then \( 49S = 0 \) is automatic. Example: \( a_1 = \cdots = a_{51} = 1 \) over \( \mathbb{F}_7 \), giving \( b_i = 50 = 1 = a_i \). \checkmark

\textit{Case 3:} Characteristic is 2. Then \( -a_i = a_i \), so \( b_i = S + a_i = S - a_i \). Take any 51 non-zero elements with \( S = 0 \); then \( b_i = a_i \). \checkmark \hfill\(\square\)
\\

\textbf{Problem 9.} (IMC 2006, Day 1, Problem 2)
Find the number of positive integers \( x \) satisfying:
\begin{enumerate}
    \item \( x < 10^{2006} \);
    \item \( x^2 - x \) is divisible by \( 10^{2006} \).
\end{enumerate}

\textbf{Solution.}
\textit{Answer:} 3.

We need \( 10^{2006} \mid x(x-1) \). Since \( \gcd(x, x-1) = 1 \) and \( 10^{2006} = 2^{2006} \cdot 5^{2006} \), each prime power must divide exactly one of \( x \) or \( x - 1 \).

By the Chinese Remainder Theorem, we solve four systems:
\begin{enumerate}
    \item \( x \equiv 0 \pmod{2^{2006}} \) and \( x \equiv 0 \pmod{5^{2006}} \): gives \( x \equiv 0 \pmod{10^{2006}} \), so \( x = 0 \) (excluded since \( x > 0 \)).
    \item \( x \equiv 1 \pmod{2^{2006}} \) and \( x \equiv 1 \pmod{5^{2006}} \): gives \( x \equiv 1 \pmod{10^{2006}} \), so \( x = 1 \).
    \item \( x \equiv 0 \pmod{2^{2006}} \) and \( x \equiv 1 \pmod{5^{2006}} \): by CRT, a unique \( x \) in \( \{1, \ldots, 10^{2006} - 1\} \).
    \item \( x \equiv 1 \pmod{2^{2006}} \) and \( x \equiv 0 \pmod{5^{2006}} \): by CRT, a unique \( x \) in \( \{1, \ldots, 10^{2006} - 1\} \).
\end{enumerate}

Cases 2, 3, 4 each yield one solution in \( \{1, \ldots, 10^{2006} - 1\} \), giving \textbf{3} solutions. \hfill\(\square\)
\\

\textbf{Problem 10.} (IMC 2006, Day 2, Problem 5)
Prove that there exists an infinite number of relatively prime pairs \( (m, n) \) of positive integers such that the equation
\[
(x + m)^3 = nx
\]
has three distinct integer roots.

\textbf{Solution.}
Substituting \( y = x + m \) gives \( y^3 = n(y - m) \), i.e., \( y^3 - ny + mn = 0 \). By Vieta's formulas, the three roots \( y_1, y_2, y_3 \) satisfy:
\[
y_1 + y_2 + y_3 = 0, \quad y_1 y_2 + y_1 y_3 + y_2 y_3 = -n, \quad y_1 y_2 y_3 = -mn.
\]

Set \( y_1 = p(p^2 + pq + q^2) \), \( y_2 = q(p^2 + pq + q^2) \), \( y_3 = -(p+q)(p^2 + pq + q^2) \) for positive integers \( p, q \). Then \( y_1 + y_2 + y_3 = 0 \). Let \( k = p^2 + pq + q^2 \). Then:
\[
n = -(y_1 y_2 + y_1 y_3 + y_2 y_3) = k^2(p^2 + pq + q^2) = k^3,
\]
\[
m = -\frac{y_1 y_2 y_3}{n} = \frac{pq(p+q) \cdot k^3}{k^3} = pq(p+q).
\]

The three roots in \( x \) are \( x_i = y_i - m \), which are distinct since \( y_i \) are distinct (as \( p, q > 0 \)).

It remains to show \( \gcd(m, n) = 1 \) for infinitely many \( (p, q) \). Take \( p = 1 \). Then \( m = q(q+1) \) and \( n = (1 + q + q^2)^3 \). A prime dividing both \( m \) and \( n \) divides \( q(q+1) \) and \( q^2 + q + 1 \). But \( q^2 + q + 1 = q(q+1) + 1 \), so \( \gcd(q(q+1), q^2+q+1) = 1 \). Hence \( \gcd(m, n) = 1 \) for all \( q \geq 1 \), giving infinitely many coprime pairs. \hfill\(\square\)

\end{document}
